\subsection{代数张量积的泛性质刻画}

命题 5. 设 \((\mathbf{E}_i)_{i\in I}\) 是一个有限的 A-代数族，且对每个 \(i\in I\) ，设 \(e_i\) 是 \(\mathbf{E}_i\) 的单位元。对每个 \(i\in I\) ，设 \(u_i: \mathbf{E}_i \to \mathbf{E} = \bigotimes_{i\in I} \mathbf{E}_i\) 是由下式定义的 A-线性映射：

\[
u _ {i} (x _ {i}) = \bigotimes_ {j} x _ {j} ^ {\prime} \quad \text { with } x _ {i} ^ {\prime} = x _ {i} \text { and } x _ {j} ^ {\prime} = e _ {j} \text { for } j \neq i.
\]

\begin{enumerate}
\def\labelenumi{(\roman{enumi})}
\item
  诸 \(u_{i}\) 是 A-代数同构；进而，对 \(i \neq j\) ，元素 \(u_{i}(x_{i})\) 与 \(u_{j}(x_{j})\) 在 \(\mathbf{E}\) 中对诸 \(x_{i} \in \mathbf{E}_{i}\) 与 \(x_{j} \in \mathbf{E}_{j}\) 可交换，且 \(\mathbf{E}\) 由诸子代数 \(u_{i}(\mathbf{E}_{i})\) 的并生成 .
\item
  设 \(\mathbf{F}\) 是一个 A-代数，且对所有 \(i \in \mathbf{I}\) ，设 \(v_i: \mathbf{E}_i \to \mathbf{F}\) 是一个 A-代数同态，其中诸 \(v_i\) 满足：对 \(i \neq j\) ， \(v_i(x_i)\) 与 \(v_j(x_j)\) 在 \(\mathbf{F}\) 中对诸 \(x_i \in \mathbf{E}_i\) 与 \(x_j \in \mathbf{E}_j\) 可交换。则存在且仅存在一个 A-代数同态 \(w: \mathbf{E} \to \mathbf{F}\) 使得
\end{enumerate}

\[
v _ {i} = w \circ u _ {i} \quad \text { for all } i \in I.\tag{4}
\]

\begin{enumerate}
\def\labelenumi{(\roman{enumi})}
\tightlist
\item
  映射 \(u_{i}\) 由 \(\mathbf{E}\) 上乘法的定义可知是代数同态。若 \(i \neq j, x_{i} \in \mathbf{E}_{i}, x_{j} \in \mathbf{E}_{j}\) ，则
\end{enumerate}

\[
u _ {i} (x _ {i}) = \bigotimes_ {k} x _ {k} ^ {\prime} \quad \text { with } x _ {i} ^ {\prime} = x _ {i}, x _ {k} ^ {\prime} = e _ {k} \text { for } k \neq i,
\]

\[
u _ {j} (x _ {j}) = \bigotimes_ {k} x _ {k} ^ {\prime \prime} \quad \text { with } x _ {j} ^ {\prime \prime} = x _ {j}, x _ {k} ^ {\prime \prime} = e _ {k} \text { for } k \neq j.
\]

显然 \(x_{k}^{\prime}x_{k}^{\prime \prime} = x_{k}^{\prime \prime}x_{k}^{\prime}\) 对所有 \(k\in \mathbf{I}\) 成立，因此由定义 E 中乘法的公式 (1)（第 1 号）， \(u_{i}(x_{i})\) 与 \(u_{j}(x_{j})\) 在 E 中交换。最后一个断言由关系式 \(\bigotimes_{i}x_{i} = \prod_{i\in \mathbf{I}}u_{i}(x_{i})\) 得出 .

\begin{enumerate}
\def\labelenumi{(\roman{enumi})}
\setcounter{enumi}{1}
\tightlist
\item
  对每个 \(i \in I\) ，设 \(x_{i}\) 是 \(\mathbf{E}_{i}\) 的一个元素。乘积 \(\prod_{i \in I} v_{i}(x_{i})\) 在 \(\mathbf{F}\) 中的定义与 \(\mathbf{I}\) 上的任何次序无关，因为代数 \(\mathbf{F}\) 是结合的，且诸元素 \(v_{i}(x_{i})\) 两两可交换。映射 \((x_{i})_{i \in I} \to \prod_{i \in I} v_{i}(x_{i})\) （从 \(\prod_{i \in I} \mathbf{E}_{i}\) 到 \(\mathbf{F}\) ）显然是 A-多重线性的，因此存在且仅存在一个 A-线性映射 \(w: \mathbf{E} \to \mathbf{F}\) 使得
\end{enumerate}

\[
w \left(\bigotimes_ {i} x _ {i}\right) = \prod_ {i} v _ {i} \left(x _ {i}\right).\tag{5}
\]

现在，所求的 A-代数同态 \(w: \mathbf{E} \to \mathbf{F}\) 必须满足 (5)，这由 (4) 及事实 \(\bigotimes_{i} x_i = \prod_{i \in I} u_i(x_i)\) 得出。这就证明了 \(w\) 的唯一性；还需证明由 (5) 定义的 A-线性映射 \(w\) 是 A-代数同态且满足 (4)。 \(w\) 满足 (4) 是显然的：只需把 (5) 应用于 \(x_j = e_j\) （ \(j \neq i\) ）的情形，就得到 \(w(u_i(x_i)) = v_i(x_i)\) 。最后， \(w\) 是代数同态，因为

\[
\begin{array}{r l} w \left(\left(\bigotimes_ {i} x _ {i}\right) \left(\bigotimes_ {i} y _ {i}\right)\right) & = w \left(\bigotimes_ {i} (x _ {i} y _ {i})\right) = \prod_ {i} v _ {i} (x _ {i} y _ {i}) \\ & = \prod_ {i} (v _ {i} (x _ {i}) v _ {i} (y _ {i})) = \left(\prod_ {i} v _ {i} (x _ {i})\right). \left(\prod_ {i} v _ {i} (y _ {i})\right) \end{array}
\]

因为 \(v_{i}(x_{i})\) 与 \(v_{j}(y_{j})\) 对 \(j \neq i\) 交换；从而

\[
w \left(\left(\bigotimes_ {i} x _ {i}\right) \left(\bigotimes_ {i} y _ {i}\right)\right) = w \left(\bigotimes_ {i} x _ {i}\right). w \left(\bigotimes_ {i} y _ {i}\right)
\]

由线性性，这就完成了证明。

由 E 与典范映射 \(\phi: (x_i) \mapsto \bigotimes_i x_i\) （从 \(\prod_i E_i\) 到 E 的映射）组成的有序对，是泛映射问题（集合论, IV, §3, 第 1 号）的一个解，其中 \(\Sigma\) 是 A-代数结构的种类，态射是 A-代数同态，而 \(\alpha\) -映射是诸映射 \(\prod_i u_i\) （从 \(\prod_i E_i\) 到某个 A-代数的映射），它们使得诸 \(u_i\) 是 A-代数同态，且 \(u_i(x_i)\) 与 \(u_j(x_j)\) 对 \(i \neq j\) 、对所有 \(x_i \in E_i\) 和 \(x_j \in E_j\) 交换 .

推论. 设 \((\mathbf{E}_i)_{i\in I}, (\mathbf{F}_i)_{i\in I}\) 是 A-代数的两个有限族，且对所有 \(i\in I\) ，设 \(f_i: \mathbf{E}_i \to \mathbf{F}_i\) 是代数同态。若 \(u_i: \mathbf{E}_i \to \bigotimes_{j\in I} \mathbf{E}_{j}, v_i: \mathbf{F}_i \to \bigotimes_{j\in I} \mathbf{F}_j\) 是典范同态，则映射 \(f = \bigotimes_{i} f_i\) （参见第 1 号）是唯一的 A-代数同态，使得 \(f \circ u_i = v_i \circ f_i\) 对所有 \(i\in I\) 成立 .

只需注意，同态 \(g_{i} = v_{i} \circ f_{i}\) 满足 \(g_{i}(x_{i}) = v_{i}(f_{i}(x_{i}))\) 与 \(g_{j}(x_{j}) = v_{j}(f_{j}(x_{j}))\) 对 \(i \neq j, x_{i} \in \mathbf{E}_{i}\) 和 \(x_{j} \in \mathbf{E}_{j}\) 可交换；然后应用命题 5。

当在命题 5 中假设代数 F 是交换的时， \(v_{i}(x_{i})\) 与 \(v_{j}(x_{j})\) 对 \(i \neq j\) 可交换这一假设自动满足。因此，当 F 交换时，存在典范双射

\[
\operatorname{Hom} _ {\mathrm{A-alg.}} \left(\bigotimes_ {i} \mathrm{E} _ {i}, \mathrm{F}\right)\rightarrow \prod_ {i} \operatorname{Hom} _ {\mathrm{A-alg.}} (\mathrm{E} _ {i}, \mathrm{F}),\tag{6}
\]

即把每个同态 \(w\) （ \(\bigotimes_{i} E_{i}\) 到 \(F\) 的）对应到诸 \(w \circ u_{i}\) 所成之族的那个双射 .

注意，若 \(\mathbf{E}\) 是交换 A-代数，则 \(\mathbf{E} \otimes_{\mathbf{A}} \mathbf{F}\) 的环结构与 \(\mathbf{F}_{(\mathbf{E})}\) 的环结构相同（ \(\S 1\) , 第 5 号）。

